\couple{Anche se fosse}{Un giorno un ladro chiese ad un mercante se la sua frutta fosse dolce, ed il mercante disse “sì”. Quello stesso giorno il ladro mercante chiese al mercante ladro se la sua frutta fosse dolce, e il mercante ladro disse “no”. Poche ore dopo un cliente chiese al mercante ladro se la sua frutta fosse dolce, ed a quel punto il mercante ladro disse “sì”. Tale cliente era a conoscenza dell’astuzia dei mercanti e della malizia dei ladri, perciò continuò chiedendo se tale frutta venisse da lontano.
A quel punto il ladro mercante disse “no sì”, mentre il mercante ladro disse “sì no”. Il cliente, sentendo l’odore di menzogna, decise di intascarsi un paio di mele di nascosto e partire senza pagare. Arrivò a casa e le posò alla finestra. Una colomba chiese al cliente se le sue mele fossero dolci ed il cliente disse “sì”. Il merlo chiese alla colomba se la mela che aveva rubato al cliente fosse dolce e la colomba rispose “no”. L’altra mela venne conquistata dalla muffa prima che il cliente potesse assaggiarla, come poteva conoscerne il gusto? Questi clienti, ladri e mercanti tutti a parlare di nulla solo per potersi mentire. Forse il tale aveva chiesto alla muffa se la mela fosse dolce e la muffa gli aveva risposto “fatti i cazzi tuoi”. E fu così che, presa dall’azione, la muffa chiese a te se tutta sta frutta, che a questo punto forse non è frutta , fosse dolce oppure no, e tu non dicesti altro che “non ho capito un cazzo”.}{One day a thief asked a merchant if his fruit was sweet, and the merchant said “yes.” That same day, the thief merchant asked the merchant thief if his fruit was sweet, and the merchant thief said “no.” A few hours later, a customer asked the merchant thief if his fruit was sweet, and this time the merchant thief said “yes.” The customer knew well the cunning of merchants and the malice of thieves, so he went on to ask whether the fruit came from afar. At that point, the thief merchant said “no yes,” while the merchant thief said “yes no.” The customer, smelling the stench of lies, decided to pocket a couple of apples and leave without paying. He got home and set them on the windowsill. A dove asked the customer if his apples were sweet, and the customer said “yes.” The blackbird asked the dove if the apple it had stolen from the customer was sweet, and the dove replied “no.” The other apple was conquered by mold before the customer could taste it, how could he ever know its flavor? These customers, thieves, and merchants all talking nonsense just to lie to one another. Maybe the man had asked the mold if the apple was sweet, and the mold had answered, “mind your own fucking business.” And so it was that, carried by action, the mold turned to you and asked whether all this fruit, which by now perhaps isn’t fruit at all, was sweet or not, and you said nothing but, “I didn’t understand shit.”}

\couple{Sensei}{「師よ、真とは何ぞや。」
師、ただ肩を竦（すく）めたり。

「然らば、我いかに申すべきや。」
「此の紙に記し、しかして之に糞せよ。」

弟子、意見と糞にて穢れたる紙を奉（たてまつ）る。
師、これを恭（うやうや）しく舐め、拝し、
額に打ち釘ちて、村々を巡り歩き、
声尽き、力尽き、命尽くるまでこれを唱えたり。

最期に囁きて曰く、
「これぞ真なり。」}{"Master, what is the truth?" 
The master shrugged his shoulders.
“Then what should I say?”
"Write it in on this paper and cover it with feces."

The student offered the toilet paper dirtied by opinions and excrement. 
The master licked it devotedly, bowed to it, nailed it into his forehead.
He wandered from village to village, chanting it, 
until he lost his voice, strength and life, just after whispering "this is the Truth".}

\couple{P.}{il giorno di P. la sveglia non suono’ o perlomeno non venne sentita. forse spenta nel sonno. 

il giorno di P. la mamma chiamo’ per essere rassicurata che non ci fossero stati intoppi fino a P. 

il giorno di P. la gente scese in strada a protestare contro l’assedio di P. 

il giorno di P. l’aereo perse quota ma non slancio, e qualcuno grido’ allo scandalo.

Il giorno di P. stettero tutti in silenzio ad aspettare P.

il giorno prima di P. c’erano tensione fremito - e poco P.

il giorno dopo di P. le bici avevano le banane alle ruote e versarono P.

il giorno di P. c’e’ Nanni Moretti in prima fila che discute con P. 

dal giorno di P. abbiamo avuto solo giorni pieni di P.

Nessun giovane vuole P. e nessun vecchio vuole non-P. 

I bambini giocano con P.}{The dawn of P came without anybody noticing it.

Soon after dawn it started raining P on the roofs and canals of the urban conglomerate. 

More and more P poured and soon enough the streets were flooded with P. 

P was flowing into the basements and leaking through the roofs. 

Concrete cubic containers with effortless modern interiors were now taken over by P. 

P did not like the effortlessly-modern interiors. P preferred a rustic and traditional design.

P expanded to fill the interior and, as it retracted, the head of a stag over a Victorian fireplace appeared.    

The original dwellers thought it best not to resist any choice made by P.

P seemed all-knowledgeable and benevolent; the dwellers started to worship P.

The same events unfolded in all other concrete cubic containers. 

With time P had become indiscernible from its surroundings. It was the sunset of P. 

It will soon be the dawn of a new P.}

